%! Orthogonal Matrices and Gram--Schmidt — Unit 4, Lesson 4.4 — Student Packet Cover Sheet
\documentclass[10pt]{article}
\usepackage{linalg-article}
\usepackage{linalg-boxes}
\usepackage{ltablex}
\keepXColumns

\begin{document}

% Full-bleed burgundy banner
\begin{tikzpicture}[remember picture, overlay]
  \fill[burgundy] (current page.north west)
    rectangle ([yshift=-0.9in]current page.north east);
\end{tikzpicture}
\vspace{-0.8in}

\begin{center}
  {\color{white}\textbf{\LARGE Linear Algebra}}\\[4pt]
  {\color{blushmid}\large\textbf{Unit 4 \quad Orthogonality}}\\[6pt]
  {\color{white}\textbf{\large Lesson 4.4 \quad Orthogonal Matrices and Gram--Schmidt}}
\end{center}

\vspace{0.22in}
\namedateperiod
\vspace{0.18in}

\begin{learningtargetbox}
\small
\begin{enumerate}[leftmargin=1.5em, itemsep=5pt]
  \item \ldots{} test whether vectors are \textbf{orthonormal} (perpendicular \emph{and} unit length) and build the matrix $Q$ with $Q^{\mathsf{T}}Q=I$.
  \item \ldots{} find coordinates and projections with \textbf{no solving} --- each weight is a dot product $q_i\cdot\mathbf{b}$, so $\hat{\mathbf{x}}=Q^{\mathsf{T}}\mathbf{b}$.
  \item \ldots{} run \textbf{Gram--Schmidt} to turn independent vectors into an orthonormal set, and explain \emph{why} the subtraction step makes them perpendicular.
\end{enumerate}
\end{learningtargetbox}

\vspace{0.14in}

\begin{tocbox}
\small
\renewcommand{\arraystretch}{1.5}
\begin{tabularx}{\linewidth}{@{} c l X r @{}}
\rowcolor{blush}
\textbf{\#} & \textbf{Component} & \textbf{Description} & \textbf{Score} \\
\midrule
1 & Warm-Up        & Spiral review: length \& unit vectors, the perpendicular test, the normal equations & \blank{1.2cm} \\
2 & Guided Notes   & Orthonormal vectors, $Q^{\mathsf{T}}Q=I$, coordinates by dot product, Gram--Schmidt & \blank{1.2cm} \\
3 & Group Activity & Perfect Axes --- test orthonormality, coordinates for free, build $Q$ by Gram--Schmidt & \blank{1.2cm} \\
4 & Exit Ticket    & Quick check: orthonormality test, a normalize step, and why $Q^{\mathsf{T}}Q=I$ helps & \blank{1.2cm} \\
5 & Homework       & Normalizing, coordinates by dot product, Gram--Schmidt, and $A=QR$ & \blank{1.2cm} \\
\midrule
  & \multicolumn{2}{r}{\textbf{Total}} & \blank{1.2cm} \\
\end{tabularx}
\end{tocbox}

\vspace{0.10in}

\begin{remindbox}
\small One idea drives the lesson: \textbf{orthonormal} vectors are ``perfect graph paper'' ---
perpendicular \emph{and} length~$1$. Stacked into a matrix $Q$, they satisfy $Q^{\mathsf{T}}Q=I$, so every
coordinate is a single dot product $q_i\cdot\mathbf{b}$ and last lesson's normal equations collapse to
$\hat{\mathbf{x}}=Q^{\mathsf{T}}\mathbf{b}$ --- projection with \emph{no} solving. \textbf{Gram--Schmidt}
manufactures such vectors from any independent set by subtracting off overlaps and normalizing.
\end{remindbox}

\end{document}
